\documentclass[a4paper,11pt]{article}

\usepackage[utf8]{inputenc}
\usepackage[T1]{fontenc}
\usepackage[french]{babel}
\usepackage{caption}
\usepackage{subcaption}

\usepackage{amsmath}
\usepackage{amssymb}
\usepackage{graphicx}
\usepackage{geometry}
\usepackage{hyperref}
\usepackage{listings}
\usepackage{xcolor}
\usepackage{booktabs}
\usepackage{tikz}

\geometry{margin=2.5cm}

\lstset{
language=Python,
basicstyle=\ttfamily\small,
keywordstyle=\color{blue},
commentstyle=\color{gray},
stringstyle=\color{red},
frame=single,
breaklines=true
}

\title{\textbf{Transformation de coordonnées entre repère objet et repère robot}}
\author{\phantom{Mariéthoz Cédric}}
\date{\today}

\begin{document}

\maketitle

\begin{abstract}

Ce rapport traite de la problématique de conversion de coordonnées entre un repère local associé à un objet et le repère global utilisé par un système robotique.
Pour établir cette correspondance, une méthode d’estimation de la transformation affine est présentée. Elle s’appuie sur un ensemble de points homologues mesurés dans les deux repères et sur une résolution par moindres carrés.

Le document aborde également une procédure permettant de convertir une normale en une orientation exploitable par le robot, afin d’assurer une manipulation cohérente de l’objet dans l’espace.

\end{abstract}

\tableofcontents
\newpage

\section{Introduction}

Dans de nombreuses applications, les coordonnées d’un objet sont exprimées dans un repère propre à cet objet.  
Cependant, les mouvements du robot sont définis dans un repère global différent.

Afin de permettre l'interaction entre ces deux systèmes, il est nécessaire de déterminer une transformation permettant de convertir les coordonnées d’un repère vers l’autre.

Cette transformation permet par exemple :

\begin{itemize}
\item d’exécuter des trajectoires définies dans le repère de l’objet
\item d’utiliser des données provenant de capteurs ou de modèles 3D
\item d’adapter automatiquement la position d’un objet dans l’espace
\end{itemize}

\section{Principe général}

La transformation entre deux repères tridimensionnels peut être décrite par une transformation affine :

\[
T =
\begin{bmatrix}
R & t \\
0 & 1
\end{bmatrix}
\]

où :

\begin{itemize}
\item $R$ est la matrice de rotation
\item $t$ est le vecteur de translation
\end{itemize}

Cette matrice permet de convertir un point $p_A$ exprimé dans le repère objet en un point $p_B$ dans le repère robot :

\[
p_B = T p_A
\]

\section{Collecte des points}

Pour estimer la transformation, il est nécessaire de disposer d’un ensemble de points correspondants dans les deux repères.

Chaque point est mesuré :

\begin{itemize}
\item dans le repère objet
\item dans le repère robot
\end{itemize}

Un minimum de trois points non colinéaires est nécessaire pour déterminer une transformation.

Cependant, l’utilisation d’un nombre plus élevé de points permet d’obtenir une solution plus robuste en utilisant une résolution par moindres carrés.

Les données collectées sont donc de la forme :

\[
A_i = (x_i, y_i, z_i)
\]

\[
B_i = (x_i', y_i', z_i')
\]

\section{Calcul de la transformation}

La transformation recherchée doit satisfaire :

\[
B = RA + t
\]

\newpage
Afin de résoudre ce problème, on construit une matrice :

\[
P =
\begin{bmatrix}
x_1 & y_1 & z_1 & 1 \\
x_2 & y_2 & z_2 & 1 \\
\vdots & \vdots & \vdots & \vdots
\end{bmatrix}
\]

Le problème devient alors :

\[
P X = B
\]

où :

\[
X =
\begin{bmatrix}
R^T \\
t
\end{bmatrix}
\]

La solution est obtenue par moindres carrés :

\[
X = (P^T P)^{-1} P^T B
\]

La matrice de transformation finale est ensuite reconstruite sous la forme :

\[
T =
\begin{bmatrix}
R & t \\
0 & 1
\end{bmatrix}
\]

\section{Transformation des normales}

Lorsqu’un point possède également une normale de surface, celle-ci doit être transformée correctement.

Contrairement aux positions, les normales ne doivent pas être transformées directement par la matrice $R$.

La transformation correcte est :

\[
n' = (R^{-1})^T n
\]

Cette opération correspond à la transformation des vecteurs normaux sous une transformation affine.

\section{Conversion des normales en orientation}

Dans certaines applications, la normale d’une surface doit être convertie en orientation utilisable par le robot.

\subsection{Conversion d'une normale en angles d'Euler}

\[
\text{Soit } \mathbf{n} = (n_x, n_y, n_z) \in \mathbb{R}^3
\quad\text{une normale unitaire telle que}\quad
\|\mathbf{n}\| = 1.
\]

L'objectif est de déterminer des angles d'Euler 
\((\alpha, \beta, \gamma)\) dans la convention 
\(\text{Yaw--Pitch--Roll} = (Z,Y,X)\)
tels que la rotation composée

\[
R = R_z(\alpha)\, R_y(\beta)\, R_x(\gamma)
=
\begin{bmatrix}
\cos\alpha & -\sin\alpha & 0 \\
\sin\alpha &  \cos\alpha & 0 \\
0          &  0          & 1
\end{bmatrix}
\begin{bmatrix}
\cos\beta & 0 & \sin\beta \\
0         & 1 & 0         \\
-\sin\beta & 0 & \cos\beta
\end{bmatrix}
\begin{bmatrix}
1 & 0          & 0           \\
0 & \cos\gamma & -\sin\gamma \\
0 & \sin\gamma &  \cos\gamma
\end{bmatrix}
\]

\newpage

aligne l'axe canonique

\[
\mathbf{e}_z = (0,0,1)
\]

avec la normale \(\mathbf{n}\).

\[
R_z(\alpha)\, R_y(\beta)\,
\begin{pmatrix}
0 \\ 0 \\ 1
\end{pmatrix}
=
\begin{pmatrix}
\sin\beta \cos\alpha \\
\sin\beta \sin\alpha \\
\cos\beta
\end{pmatrix}.
\]

En identifiant ce vecteur à la normale \(\mathbf{n}\), on obtient le système :

\[
n_x = \sin\beta \cos\alpha,
\qquad
n_y = \sin\beta \sin\alpha,
\qquad
n_z = \cos\beta.
\]

On en déduit directement les angles :

\[
\beta = \arctan2\!\left(\sqrt{n_x^2 + n_y^2},\, n_z\right)
\qquad\text{(inclinaison / pitch)}
\]

\[
\alpha = \arctan2(n_y,\, n_x)
\qquad\text{(orientation horizontale / yaw)}
\]

Le troisième angle \(\gamma\) (roll) représente une rotation autour de la normale elle-même.  
Comme cette rotation n'affecte pas l'alignement de l'axe \(z\) avec \(\mathbf{n}\), on peut choisir :

\[
\gamma = 0
\]

Ainsi, la rotation d'Euler

\[
R = R_z(\alpha)\, R_y(\beta)
\]

est la transformation minimale qui aligne l'axe \(z\) avec la normale \(\mathbf{n}\).

\subsection{Conversion d'une normale en vecteur de rotation}

Une autre représentation des orientations consiste à utiliser un vecteur de rotation.

Dans cette représentation, une rotation est décrite par :

\begin{itemize}
\item un axe de rotation unitaire $k$
\item un angle de rotation $\theta$
\end{itemize}

Le vecteur de rotation est défini par :

\[
r = \theta k
\]

Afin d'aligner l'axe $Z$ avec la normale $n = (x,y,z)$, l'angle de rotation est :

\[
\theta = \arccos(z)
\]

L'axe de rotation correspond au produit vectoriel entre l'axe $Z$ et la normale :

\[
k = Z \times n
\]

avec :

\[
Z =
\begin{bmatrix}
0 \\
0 \\
1
\end{bmatrix}
\]

\newpage

Ce qui donne :

\[
k =
\begin{bmatrix}
-y \\
x \\
0
\end{bmatrix}
\]

Après normalisation :

\[
k =
\frac{1}{\sqrt{x^2+y^2}}
\begin{bmatrix}
-y \\
x \\
0
\end{bmatrix}
\]

Le vecteur de rotation final est alors :

\[
r =
\theta
\begin{bmatrix}
-\frac{y}{\sqrt{x^2+y^2}} \\
\frac{x}{\sqrt{x^2+y^2}} \\
0
\end{bmatrix}
\]

\subsection{Comparaison des méthodes}

\begin{figure}[h]
    \centering

    \begin{subfigure}{0.45\textwidth}
        \centering
        \includegraphics[width=\linewidth]{Euler_angle.png}
        \caption{Angle d'Euler}
        \footnotetext{\tiny \url{https://en.wikipedia.org/wiki/Euler_angles}}

        \label{fig:sub1}
    \end{subfigure}
    \hfill
    \begin{subfigure}{0.49\textwidth}
        \centering
        \includegraphics[width=\linewidth]{Rotation_vector.png}
        \caption{Vecteur de rotation}
        \footnotetext{\tiny \url{https://math.stackexchange.com/questions/511370/how-to-rotate-one-vector-about-another}}
        \label{fig:sub2}
    \end{subfigure}

    \caption{Représentation des différences des principe de rotation}
    \label{fig:main}
\end{figure}


Une normale définit uniquement une direction dans l'espace et ne fournit qu'une contrainte partielle sur l'orientation.

La conversion en angles d'Euler nécessite cependant de définir trois rotations indépendantes.  
Dans l'approche précédente, la rotation autour de l'axe $Z$ est fixée arbitrairement :

\[
r_z = 0
\]

\newpage

Cette hypothèse impose une orientation particulière autour de la normale, ce qui ne correspond pas nécessairement à la géométrie réelle du problème.

De plus, les angles d'Euler peuvent présenter des discontinuités lorsque certaines orientations sont atteintes.

La représentation par vecteur de rotation présente plusieurs avantages :

\begin{itemize}
\item elle décrit directement une rotation minimale entre deux directions
\item elle évite l'introduction d'hypothèses arbitraires
\item elle correspond à la représentation utilisée par le robot
\end{itemize}

Pour ces raisons, la conversion de la normale en vecteur de rotation constitue une solution plus adaptée dans ce contexte.

\section{Implémentation}

Algorithme:

\begin{enumerate}

\item Mesurer les points de référence dans l'espace de l'objet dans l'espace du robot
\item Calculer la matrice de transformation pour passer d'un espace à l'autre
\item Transformer les coordonnées de l'espace de l'objet dans celui du robot
\item Convertir la normal en vecteur de rotation

\end{enumerate}

\newpage

\section{Conclusion}

La transformation entre repère objet et repère robot constitue une étape fondamentale pour l’intégration de données provenant de différentes sources.

La méthode présentée permet d’estimer efficacement cette transformation à partir d’un ensemble de correspondances de points.

L’utilisation des vecteurs de rotation pour représenter l’orientation offre une solution robuste et compatible avec les systèmes robotiques modernes.

\appendix

\section{Code Python}

\subsection{Calcul de la matrice de transformation}
\begin{lstlisting}

def create_transformation(A, B):

    P = np.hstack([A, np.ones((N, 1))])

    X, _, _, _ = np.linalg.lstsq(P, B, rcond=None)

    R = X[:3, :].T
    t = X[3, :]

    T = np.eye(4)
    T[:3, :3] = R
    T[:3, 3] = t
    
    return T


\end{lstlisting}

\subsection{Conversion normal en vecteur de rotation}
\begin{lstlisting}

def normalize(v):
    v = np.asarray(v, dtype=float)
    n = np.linalg.norm(v)
    return v / n if n > 0 else v

def normal_to_rxyz(n):

    x, y, z = normalize(n)

    theta = np.arccos(z)
    r = np.sqrt(x*x + y*y)

    rx = -theta * (y / r)
    ry =  theta * (x / r)
    rz = 0

    return np.array([rx,ry,rz])

\end{lstlisting}

\end{document}